% D proof; root and B independent literal audits passed.
\subsection{Seven excess arrows with acyclic support}

\begin{lemma}[Two support-homology budgets]
\label{C:flag-second-cycle-budget}
Let $\Pi$ be a finite connected poset and let $R$ be the cycle
rank of its undirected comparability graph.  If $R\leq6$, then
$H_2(B\Pi;\mathbb Q)=0$.  If $R\leq8$, then
$H_4(B\Pi;\mathbb Q)=0$.
\end{lemma}
\begin{proof}
The order complex is the flag complex of the comparability graph.
Suppose it has nonzero second homology, and choose an induced
subcomplex with this property and the fewest vertices.  It is
connected.  Deleting any vertex kills its second homology.
The Mayer--Vietoris sequence for the deletion and the vertex
star therefore injects its nonzero second homology into the
first homology of the vertex link.  A flag complex with nonzero
first homology has at least four vertices: on at most three
vertices every graph cycle is a filled triangle.  Every vertex
of our induced graph consequently has degree at least four.
There are at least six vertices, since a five-vertex graph of
minimum degree four is complete and its flag complex is a simplex.
Writing $v,e$ for these vertex and edge counts gives
$e-v+1\geq v+1\geq7$.  Graph cycle rank cannot increase on
passing to a subgraph, proving the first assertion.

A four-simplex requires a complete five-vertex subgraph.
Two different such subgraphs have union of cycle rank at least
nine: the smallest value occurs when they share four vertices,
and then the union has six vertices and fourteen edges.
Disjoint complete subgraphs already contribute twelve cycles.
Thus $R\leq8$ permits at most one four-simplex and no
five-simplex.  The boundary of that one four-simplex, if present,
is nonzero, so fourth homology vanishes.
\end{proof}

\begin{lemma}[A cut through the parallel hom-sets]
\label{C:parallel-hom-cut}
Let $C$ be a finite category with acyclic support of cycle rank
at most six.  Suppose every corner has zero rational second
cohomology.  Call a nonempty cross hom-set parallel if it has
at least two arrows.  Suppose no two parallel hom-sets
$C(x,y)$ and $C(z,w)$ satisfy $y\leq z$ in the support order.
Then every product of two rational degree-two classes on $BC$
vanishes.
\end{lemma}
\begin{proof}
If there are no parallel hom-sets, apply
Lemma~\ref{C:singleton-support-H2} componentwise and
Lemma~\ref{C:flag-second-cycle-budget}.
Otherwise let $A$ be the downward closure of all sources of
parallel hom-sets, and let $D$ be its complement.  The stated
condition says precisely that every parallel hom-set crosses
from $A$ to $D$.  There are no arrows from $D$ to $A$.
Every connected component of either full block having at least
two objects has singleton cross hom-sets.  Its second
cohomology is that of its support by
Lemma~\ref{C:singleton-support-H2}, and is zero by the first
part of Lemma~\ref{C:flag-second-cycle-budget}.  A single-object
component has zero second cohomology by hypothesis.  Both
blocks therefore have zero second cohomology, and
Lemma~\ref{C:two-block-zero-H2-products} applies.

We shall also use the same proof when a corner has nonzero
second cohomology but the restriction of every global
degree-two class to that corner is zero.  Every global class
then restricts to zero on both full blocks, and the actual
two-block suspension quotient still pulls back all those
classes.  No vanishing assertion about that corner is needed.
\end{proof}

\begin{theorem}\label{C:acyclic-arrow-excess-seven}
Let $C$ have $k$ objects, $n$ arrows, and acyclic support.
If $n\leq2k+6$, then $BC\not\simeq\mathbb CP^2$.
In particular the acyclic seven-object case is excluded
through twenty arrows.
\end{theorem}
\begin{proof}
Suppose $BC\simeq\mathbb CP^2$.  We may assume that $C$ is
connected and $k\geq3$, since the cases $k\leq2$ have at
most ten arrows and are covered by the global lower bound.
Use the corner, parallel-arrow, and support excesses
\[
 D=\sum_x(|\operatorname{End}(x)|-1),\qquad
 P=\sum_{x<y}(|C(x,y)|-1),\qquad R=E-k+1.
\]
Thus $D+P+R=n-(2k-1)\leq7$.
The support functor is onto on fundamental groups, by lifting
edge loops to actual arrows and formal inverses.  Hence its
order complex has zero first homology over $\mathbb F_2$.

If $P=0$, every degree-two class is pulled back from the
support poset by Lemma~\ref{C:singleton-support-H2}.
Its fourth rational homology is zero by
Lemma~\ref{C:flag-second-cycle-budget}, so all these squares
vanish.

Suppose all corner algebras are semisimple.  If $R\leq5$,
the relative bar has strict-object length at most three,
so fourth homology vanishes.  The only remaining budget is
$D=0$, $P=1$, $R=6$.  There is at most one five-object
chain, and all corners are trivial.  A nondegenerate
four-simplex uses its four successive arrows.  If the unique
doubled hom-set is not one of these four, there is one such
simplex and its boundary is nonzero.  Otherwise there are
two such simplices.  Choose a face deleting a vertex which
is not an endpoint of the doubled successive hom-set.
That face retains its arrow label and distinguishes the two
simplices.  No other four-simplex can cancel either face,
since there is no other five-object chain.  In either case
the degree-four boundary is injective.  Thus $H_4=0$ here too.

We may now assume a nonsemisimple corner, so $D\geq2$.
If $P=1$, there is only one parallel hom-set, of size two,
and $R\leq4$.  At most one corner has nonzero second
rational cohomology, since two such corners would cost
$D\geq8$.  If none does, Lemma~\ref{C:parallel-hom-cut}
applies.  If the exceptional corner has a singleton incident
hom-set, that two-object subcategory is a join, so every
global degree-two class restricts to zero at the corner.
The final paragraph of that lemma again applies.
Otherwise the corner is a leaf on the unique doubled edge.
Its other endpoint has order at most four, hence is
rationally acyclic.  The remainder of the category has all
cross hom-sets singleton and support cycle rank at most six;
its second cohomology is zero.  The leaf edge has zero
degree-two products by
Theorem~\ref{C:two-cross-arrows-universal}.  The actual
nerve union meets at the rationally acyclic other corner,
so fourth-cohomology restriction is injective.  The global
square is therefore zero.

Henceforth $D,P\geq2$, and consequently $R\leq3$.
The support is chordal by Lemma~\ref{C:small-cycle-support}.
Its bridge and clique pieces are actual full subcategories,
whose nerves form the corresponding tree of unions.
We shall use rationally acyclic corner intersections to
inherit bounded homology and degree-two rank at most one
on a piece.  When an intersection merely has zero third
homology, fourth-cohomology restriction is still injective.
For boundedness of a piece it is enough that the intersections
have bounded rational homology.  All these statements follow
from the Mayer--Vietoris sequence of the actual nerve unions.

\paragraph{Tree support.}
If $R=0$ and every corner has zero second cohomology,
Corollary~\ref{C:height-one-zero-H2-products} applies.
There can otherwise be only one corner $M$ with nonzero
second cohomology.  A singleton incident edge kills every
global restriction to that corner, as before; the
height-one suspension quotient then kills all global products.
We may assume every incident edge has at least two arrows.
Since $P\geq2$ and $D+P\leq7$, this corner has order at
most six and every other corner has order at most two.
All its incident arrows have the same direction.

The other pieces of the support tree have semisimple
corners, one strict crossing, and no homology in degrees
at least two.  They meet the central star at rationally
acyclic corners.  The star thus inherits bounded rational
homology and degree-two rank at most one.  Exact coinvariants
at its order-at-most-two leaves give an actual rational
cohomology-ring equivalence to a star with trivial leaves.
Identify these leaf points.  The result is an ordinary
act cone over $M$; on realization the identification adds
only circle summands.  It preserves cohomology and products
in degrees at least two and preserves boundedness.
Lemma~\ref{C:nine-bounded-rank-one}, with $|M|\leq6$,
kills the square.  Fourth-cohomology restriction across
the tree then kills the global square.

\paragraph{One independent cycle.}
Let $R=1$.  Here $D+P\leq6$, so $D\leq4$ and
every corner has order at most five.  Such corners have
zero third rational homology.  The only non-bridge block
is a triangle, and fourth-cohomology restriction to it
and the bridges is injective.

A bridge whose endpoint corners have zero second
cohomology has zero degree-two products by the two-block
lemma, regardless of its hom-set size.  If an endpoint
has nonzero second cohomology, it is the five-element
identity-adjoined $2\times2$ band, every other corner is
trivial, and the bridge is an ordinary act cone over this
two-row monoid.  Its products vanish by
Theorem~\ref{C:two-row-act-cups}.

If $D\leq3$, a triangle has at most one corner of order
at least three, of order at most four.  At an endpoint,
cut this zero-$H^2$ corner from the opposite semisimple
two-object block.  In the middle, use inherited boundedness
and Corollary~\ref{C:semisimple-ends-middle-four}.
These give zero products or zero fourth homology.

Suppose $D=4$, so $P\leq2$.  If there is a five-element
corner, all others are trivial.  At the middle, a singleton
adjacent hom-set deletes that middle object; otherwise both
adjacent sets have size two and the long set is a singleton.
Triangle compression then gives a monoid with a $2\times2$
minimal band and zero degree-two products.
At an endpoint with zero $H^2$, the opposite two-object
block is semisimple and the two-block lemma applies.
At an endpoint with nonzero $H^2$, the monoid is the
five-element two-row band monoid.  If its adjacent hom-set
is a singleton, that adjacent join and the other endpoint
give a zero-$H^2$ cut.  If both adjacent sets are parallel,
they have size two and the long set is a singleton; the
trivial extreme object makes the triangle an objectwise
one- or two-state act cone, covered by
Lemma~\ref{C:two-states-objectwise-act-cups}.
In the other cases the opposite adjacent set is a singleton,
and its trivial endpoint permits the actual middle-object
deletion.  The resulting ordinary act cone has zero products
by Theorem~\ref{C:two-row-act-cups}.

It remains that all triangle corners have order at most
four.  If an adjacent hom-set is a singleton, its join
block and the opposite zero-$H^2$ corner give the cut.
Otherwise the two adjacent sets have size two and the
long set has size one.  A trivial endpoint makes the
triangle an objectwise one- or two-state categorical act
cone, so Lemma~\ref{C:two-states-objectwise-act-cups}
applies.  The same conclusion holds after exact rational
endpoint coinvariants if either endpoint has order two.
The only remaining corner orders are $(3,1,3)$.
Both trivial modules of an order-three monoid have
projective dimension at most one.  If either relevant
endpoint trivial module is projective, the strict-object
bar has total degree at most three and $H_4=0$.
Otherwise the source is the identity-adjoined two-element
left-zero band and the target its right-zero opposite.
On a two-element adjacent act the graph model in
Lemma~\ref{C:right-zero-action-graph} has nonzero first
homology only for the common constant retraction, when
it is the one-edge circle.  If either adjacent graph has
zero first homology, the triangle filtration has no total
degree-four class.  If both are circles, the category is
exactly the family of
Theorem~\ref{C:affine-chart-family-zero}.
Thus its products vanish in every case.

\paragraph{Two independent cycles.}
Let $R=2$.  Now $D+P\leq5$ and $D\leq3$.
Every corner is rationally acyclic.  Bridge and triangle
blocks are treated as above.  If $D=2$, the only nontrivial
corner has order three.  Every edge intersection then has
zero homology in degrees at least three, by the projective-
dimension-one bound for its trivial modules.  Restriction
in degree four to the triangle pieces is injective, and
each piece inherits bounded high homology.  An endpoint
order-three corner gives total bar degree at most three;
a middle order-three corner gives zero fourth homology by
Corollary~\ref{C:semisimple-ends-middle-four}.  This closes
the $D=2$ case, including its possible third parallel-arrow
excess.

We may therefore assume $D=3$, so $P\leq2$.
Consider a diamond with
shared edge $u<v$ and exclusive vertices $a,b$.
As in the excess-six proof, an edge can have nonzero
$H^2$ only if it is parallel and has an endpoint of
order at least three.  There is at most one such larger
corner in the whole category.

If $a,b<u<v$, cut $\{a,b\}$ from $\{u,v\}$.
Unless the shared edge is potentially nonzero, both
blocks have zero $H^2$.  If it is potentially nonzero,
it uses one of the at most two parallel-arrow excesses.
At least one of $a,b$ therefore has singleton hom-sets
to both $u,v$.  Its full triangle is the join of its
corner space with the shared-edge space.  This join has
zero $H^2$, even if that exclusive corner is not trivial.
Every global degree-two class hence restricts to zero
on the shared edge, and the same two-block suspension
argument applies.  The case $u<v<a,b$ is dual.

If $u<a,b<v$, a zero-$H^2$ star on either side gives
a cut from the opposite endpoint.  Avoiding both cuts
requires a potentially nonzero edge on each side.
They must meet at the unique larger exclusive corner,
say $a$.  Both its adjacent hom-sets have size two;
all other cross hom-sets in the diamond are singleton.
The corner at $a$ has order three or four, and at most
one other corner has order two.  In particular the
shared edge and the other triangle have contractible
classifying spaces: all their cross sets are singleton
and at least one of their endpoint corner spaces is a
point.  The full diamond is homotopy equivalent to the
triangle on $u,a,v$.  This triangle has semisimple
endpoints, middle order at most four, and inherited
bounded homology.  Its fourth homology is zero.

\paragraph{Three independent cycles.}
Finally $R=3$ forces $D=2$, $P=2$: one order-three
corner and all other corners trivial.  Every edge
intersection has zero homology in degrees at least
three, because the relevant trivial module of that
corner has projective dimension at most one.
Hence fourth-cohomology restriction to the clique
pieces is injective, and each piece inherits bounded
homology in sufficiently high degrees.

For triangle pieces an endpoint order-three corner
gives total bar degree at most three.  A middle
order-three corner gives zero fourth homology by
Corollary~\ref{C:semisimple-ends-middle-four} and the
inherited boundedness.  A triangle with trivial
corners is two-dimensional.
For a complete four-object block, if the only parallel
hom-set has size three, Lemma~\ref{C:parallel-hom-cut}
applies to that block.  Otherwise every cross hom-set
has at most two arrows.  At least one of its two extreme
corners is trivial.  It is therefore an objectwise
one- or two-state act cone over the other three objects,
or the opposite of such a cone.  Its products vanish
by Lemma~\ref{C:two-states-objectwise-act-cups}.
If the order-three corner lies outside the block, its
nerve has dimension at most three and $H_4=0$ directly.
Restriction now kills the global square.

Every budget case gives a zero rational square or zero
fourth rational homology, contradicting the projective
plane.  This proves the theorem.
\end{proof}
